\chapter*{符号列表}
\chaptermark{符号列表}

\section*{状态变量}
\nomenclatureitem{\textbf{符号}}{\textbf{含义}}
\nomenclatureitem{$S$}{标的资产（股票）当前价格}
\nomenclatureitem{$V(S,t)$}{欧式期权价值（BSM/CEV模型）}
\nomenclatureitem{$V(S,v,t)$}{欧式期权价值（Heston模型）}
\nomenclatureitem{$v$}{瞬时方差（Heston模型中的随机波动率状态变量）}
\nomenclatureitem{$t$}{当前时刻}
\nomenclatureitem{$T$}{期权到期时刻}
\nomenclatureitem{$\tau = T - t$}{剩余到期时间}

\section*{期权合约参数}
\nomenclatureitem{\textbf{符号}}{\textbf{含义}}
\nomenclatureitem{$K$}{期权执行价（Strike Price）}
\nomenclatureitem{$r$}{无风险利率（连续复利）}

\section*{BSM模型参数}
\nomenclatureitem{\textbf{符号}}{\textbf{含义}}
\nomenclatureitem{$\sigma$}{标的资产的常数波动率（BSM）或尺度参数（CEV）}
\nomenclatureitem{$\mu$}{标的资产的漂移率}
\nomenclatureitem{$W_t$}{标准布朗运动}

\section*{CEV模型参数}
\nomenclatureitem{\textbf{符号}}{\textbf{含义}}
\nomenclatureitem{$\beta$}{弹性参数，控制波动率对股价的依赖程度；$\beta=1$退化为BSM，$\beta<1$产生杠杆效应}

\section*{Heston模型参数}
\nomenclatureitem{\textbf{符号}}{\textbf{含义}}
\nomenclatureitem{$\kappa$}{方差过程的均值回归速度，$\kappa>0$}
\nomenclatureitem{$\theta$}{方差过程的长期均值（长期方差）}
\nomenclatureitem{$\xi$}{方差过程的波动率（vol-of-vol）}
\nomenclatureitem{$\rho$}{股价布朗运动$W_t^S$与方差布朗运动$W_t^v$的相关系数，$\rho\in(-1,1)$}
\nomenclatureitem{$v_0$}{初始瞬时方差}
\nomenclatureitem{$W_t^S$}{驱动股价过程的标准布朗运动}
\nomenclatureitem{$W_t^v$}{驱动方差过程的标准布朗运动}

\section*{多模型PINN方法符号}
\nomenclatureitem{\textbf{符号}}{\textbf{含义}}
\nomenclatureitem{$\boldsymbol{\lambda}$}{模型参数向量$(\sigma,\beta,\kappa,\theta,\xi,\rho)$}
\nomenclatureitem{$\mathcal{F}[\cdot]$}{统一PDE算子}
\nomenclatureitem{$\text{mask}$}{软掩码$\tanh(\xi/0.05)^2$，用于在BSM/CEV与Heston之间连续插值}
\nomenclatureitem{$\sigma_\text{eff}$}{等效波动率，用于加法参数化中的Black-Scholes基线计算}
\nomenclatureitem{$V_\text{BS}$}{Black-Scholes解析解，作为加法参数化的基线}
\nomenclatureitem{$\hat{V}$}{统一PINN的期权价格预测值}
\nomenclatureitem{$\text{net}(x)$}{神经网络输出的修正量（以$K$为单位）}
\nomenclatureitem{$\mathcal{L}$}{总损失函数}
\nomenclatureitem{$\mathcal{L}_\text{pde}$}{PDE残差损失}
\nomenclatureitem{$\mathcal{L}_\text{bc}$}{边界条件损失}
\nomenclatureitem{$\mathcal{L}_\text{data}$}{数据驱动锚点损失}
\nomenclatureitem{$w_\text{pde}$}{PDE残差损失的权重系数}
\nomenclatureitem{$w_\text{bc}$}{边界条件损失的权重系数}
\nomenclatureitem{$w_\text{data}$}{数据锚点损失的权重系数}
\nomenclatureitem{$N_c$}{每训练步的PDE配点数}
\nomenclatureitem{$N_b$}{边界条件采样点数}
\nomenclatureitem{$N_d$}{数据锚点数}
\nomenclatureitem{$\Phi(\cdot)$}{标准正态累积分布函数}

\section*{评估指标}
\nomenclatureitem{\textbf{符号}}{\textbf{含义}}
\nomenclatureitem{MAE}{平均绝对误差（Mean Absolute Error）}
\nomenclatureitem{RMSE}{均方根误差（Root Mean Square Error）}
\nomenclatureitem{RelMAE$_\text{ATM}$}{平值附近区域的平均相对误差}

\section*{缩写}
\nomenclatureitem{\textbf{缩写}}{\textbf{全称}}
\nomenclatureitem{BSM}{Black-Scholes-Merton模型}
\nomenclatureitem{CEV}{Constant Elasticity of Variance，常弹性方差模型}
\nomenclatureitem{PINN}{Physics-Informed Neural Network，物理信息神经网络}
\nomenclatureitem{ICPINN}{Improved Constrained PINN，改进约束物理信息神经网络}
\nomenclatureitem{LLM}{Large Language Model，大语言模型}
\nomenclatureitem{PDE}{Partial Differential Equation，偏微分方程}
\nomenclatureitem{ATM}{At-The-Money，平值期权（$S\approx K$）}
\nomenclatureitem{OTM}{Out-of-The-Money，虚值期权（看涨期权$S<K$）}
\nomenclatureitem{ITM}{In-The-Money，实值期权（看涨期权$S>K$）}
\nomenclatureitem{CIR}{Cox-Ingersoll-Ross过程，用于Heston模型的方差过程}
\nomenclatureitem{FFT}{Fast Fourier Transform，快速傅里叶变换}
\nomenclatureitem{API}{Application Programming Interface，应用程序接口}
